\section*{v2.0.0 Structural Closure Ledger and Expansion Discipline Core}
\addcontentsline{toc}{section}{v2.0.0 Structural Closure Ledger and Expansion Discipline Core}
\begin{quote}
\textbf{Normalization note.} This block is retained as a historical ledger/governance precursor. In the normalized architecture its active role belongs chiefly to the projection/translation layer and to the later formal unfolding phases that carry the actual proof-bearing closure and transport statements.
\end{quote}

This release upgrades the monolith from a proof-memory stable object to a more
explicit \emph{closure ledger}. The point is not merely to preserve ancestry in
principle, but to keep a disciplined account of how every compressed statement,
interface row, or projection sentence remains licensed by the same upstream
closure architecture. In this sense v2.0.0 turns recoverability into an
operational bookkeeping rule: admissible content should be traceable, expandable,
and relocatable without silently changing its structural status.

\subsection*{1. Closure ledger principle}
A closure ledger is the ordered record of the transports that make a statement
usable. It does not store every intermediate calculation in full; rather, it
stores the \emph{licensed path class} by which the content travels from local
structure to admissible readout. The relevant chain is now treated as a ledgered
sequence:
\[
\begin{aligned}
&\text{local triolic structure} \Longrightarrow \text{ground condition} \Longrightarrow \text{HC compensation}\\
&\Longrightarrow \text{closure confinement} \Longrightarrow \text{probe/interface filtration}\\
&\Longrightarrow \text{lemma or tool compression} \Longrightarrow \text{bounded projection/readout}.
\end{aligned}
\]
A later re-expansion is admissible only if it re-enters the same ledger class,
not merely a similar verbal narrative.

\paragraph{Ledger Proposition A (transport accountancy).}
A compressed claim is structurally mature only if one can specify the closure
ledger that licenses its transport. If the path cannot be stated, the claim may
still be suggestive, but it has not yet acquired full structural status.

\paragraph{Proof sketch.}
The earlier versions already required closure ancestry and re-expandability. The
present strengthening adds explicit accountancy: the transport steps must be
stable enough to be named, tracked, and checked. This prevents hidden leaps from
entering the proof line under the cover of familiar notation.

\subsection*{2. Expansion discipline as a release criterion}
The next layer of consolidation concerns how new material may be added. A later
expansion should no longer be judged only by local correctness or plausibility;
it should also improve the ledger quality of the monolith. Thus expansion is now
classified by whether it strengthens one of the major controlled passages:
closure to probe, probe to interface, interface to lemma, lemma to readout, or
readout back to closure ancestry.

\paragraph{Ledger Proposition B (expansion must improve auditability).}
A new derivation is preferable to an older one when it makes the same structural
content easier to audit across compression levels. Additions that only increase
volume without sharpening one of the ledger passages remain secondary.

\paragraph{Operational reading.}
This gives the document a stricter maturity test. Future versions should not add
formal bulk merely because it is available; they should add material that makes
more of the proof architecture checkable under controlled compression and
re-expansion.

\subsection*{3. Probe-ledger compatibility}
The probe channel is now interpreted as a \emph{ledger gate}. It does not simply
filter modes before projection; it also determines which later compressed forms
inherit admissibility in a controllable way. A mode that survives probe
filtration under orthogonality and HC should arrive downstream with a stable
ledger signature that can be recognized in lemma, table, and readout form.

\paragraph{Ledger Proposition C (probe selection induces ledger signature).}
If a mode is selected by the admissible probe channel, then any later compressed
representation of that mode should preserve a recognizable signature of that
selection. If the signature disappears completely, the downstream statement must
be marked as heuristic or only partially normalized.

\subsection*{4. Table and appendix discipline}
The appendix and its tables now receive a slightly sharper role again. They are
not only memory lattices; they function as \emph{ledger surfaces}. Each row,
entry, or named tool should indicate which closure path class it abbreviates and
which re-expansion channel it expects. This gives the appendix a more active
status inside the proof architecture and reduces the danger that tables drift
into mere summaries.

\paragraph{Ledger Proposition D (rows abbreviate path classes).}
A normalized table row is admissible only if it abbreviates a stable path class
rather than an isolated sentence. The row then becomes reversible in two senses:
it can be expanded upward into lemma ancestry and laterally into equivalent tool
or interface forms.

\subsection*{5. Transition to the next mathematical deepening}
With v2.0.0 the project crosses a useful threshold. The document now has a
more explicit rule for how future mathematical elaboration should be inserted:
not as detached calculations, but as ledger-preserving thickening of already
licensed paths. This is especially relevant for the larger companion roadmap,
where many later lemmas and proof tools must remain structurally synchronized.

\paragraph{Working rule for subsequent versions.}
Whenever a new proof block is added, it should be possible to say which closure
ledger it refines, which compressed interfaces it improves, and how its content
returns to the appendix without losing admissibility. Material that satisfies
all three demands strengthens the monolith in the intended structure-first way.


\clearpage
